% =====================================================================
% BAB I PENDAHULUAN
% =====================================================================

\chapter{PENDAHULUAN}
\label{ch:pendahuluan}

% Bab ini memaparkan landasan dan arah penelitian. Subbab I.1 menguraikan latar belakang permasalahan beserta celah penelitian yang memotivasi tugas akhir ini. Subbab I.2 merumuskan pertanyaan penelitian. Subbab I.3 menetapkan tujuan dan ukuran keberhasilan pencapaian. Subbab I.4 menetapkan batasan penelitian. Subbab I.5 menguraikan metodologi pelaksanaan. Subbab I.6 menyampaikan sistematika pembahasan laporan.

\section{Latar Belakang}
\label{sec:latar-belakang}

Mahasiswa berada pada fase transisi menuju kehidupan dewasa, ketika penyesuaian terhadap lingkungan sosial baru, tuntutan akademik yang meningkat, dan ketidakpastian karier berlangsung hampir bersamaan. Tekanan psikologis yang menyertai fase ini sering dianggap bagian wajar dari proses pendewasaan, tetapi data menunjukkan skala persoalan yang lebih serius. \textcite{astutik2020} melaporkan bahwa sekitar seperempat mahasiswa Indonesia mengalami gejala depresi yang secara klinis relevan dan lebih dari separuh menunjukkan gejala kecemasan, sementara \textcite{putri2024} menunjukkan beban gejala afektif ini tetap tinggi setelah pandemi berakhir. \textcite{ashilah2024} bahkan menemukan tingkat kesepian pada mahasiswa Indonesia lebih tinggi dibandingkan kelompok pembanding di Malaysia, sejalan dengan temuan \textcite{hawkleyCacioppo2010} bahwa kesepian lebih ditentukan oleh kebermaknaan hubungan yang dirasakan daripada jumlah interaksi sosial. Itulah sebabnya seorang mahasiswa dapat tetap merasa sendirian meski berada di lingkungan kampus yang ramai secara sosial.

Masalah semacam ini tidak selalu memerlukan intervensi klinis. Sebagian besar mahasiswa lebih membutuhkan ruang untuk mengekspresikan perasaan dan merefleksikan pengalaman sebelum kondisinya memburuk, dukungan yang selama ini biasa didapat dari keluarga, teman sebaya, konselor, atau komunitas, dan yang terbukti berkorelasi dengan kesejahteraan afektif yang lebih baik \parencite{pinho2020}. Namun akses terhadap dukungan profesional pun tidak selalu mudah dijangkau. \textcite{who2022} memperkirakan sebagian besar individu dengan gangguan mental di negara berpenghasilan menengah bawah belum memperoleh layanan kesehatan mental yang memadai, dan di Indonesia keterbatasan jumlah tenaga profesional memperparah kondisi ini. \textcite{eisenberg2021} menambahkan bahwa mahasiswa kerap menunda mencari bantuan karena menganggap masalahnya belum cukup berat atau karena stigma terhadap layanan kesehatan mental. Kombinasi keduanya menciptakan celah: banyak mahasiswa yang sebenarnya butuh bicara tetapi tidak terhubung ke layanan mana pun. Celah inilah yang paling mungkin diisi oleh dukungan non-klinis yang mudah diakses dan berfungsi melengkapi, bukan menggantikan, layanan profesional.

Salah satu bentuk dukungan non-klinis yang dapat diberikan bisa berbentuk pemeriksaan suasana hati harian. Hal ini dapat memberi pemahaman tentang perasaan mahasiswa pada suatu hari, tetapi tidak dirancang untuk membaca pola gejala yang berkembang selama beberapa minggu. Padahal, gejala depresi pada mahasiswa perlu dibedakan dari perubahan suasana hati sesaat melalui instrumen yang memiliki dasar pengukuran. PHQ-9 relevan untuk kebutuhan ini karena menanyakan frekuensi sembilan gejala dalam dua minggu terakhir dan telah digunakan sebagai instrumen skrining depresi, termasuk dalam validasi berbahasa Indonesia. Karena pola gejala tidak cukup dibaca sekali, penyampaiannya perlu dapat ditawarkan kembali pada jeda yang jelas dan dilakukan secara bertahap mengikuti sembilan butir instrumen. Tantangannya bukan mengubah dukungan ini menjadi layanan diagnosis, melainkan menyampaikan skrining secara sukarela agar mahasiswa tetap memahami batas hasilnya serta tidak merasa tiba-tiba menjadi bentuk formulir klinis.

Perkembangan kecerdasan buatan generatif dalam beberapa tahun terakhir membuka salah satu kemungkinan untuk mengisi ruang tersebut, yaitu \textit{companionship chatbot} (selanjutnya disebut chatbot pendamping): sistem percakapan yang dirancang bukan untuk menjawab pertanyaan faktual, melainkan untuk menemani, mendengarkan, dan merespons kondisi emosional pengguna dalam batas non-klinis. Interaksi berulang dengan sistem semacam ini bahkan terbukti berhubungan dengan penurunan tingkat kesepian pada sebagian penggunanya \parencite{defreitasLoneliness2024}. Kajian terhadap sistem chatbot pendamping dan sistem dukungan kesehatan mental digital menunjukkan bahwa beberapa kemampuan penting sudah muncul secara terpisah: latihan reflektif berbasis CBT, validasi emosi, dukungan proaktif, skrining suasana hati secara percakapan, dan hubungan jangka panjang dengan pendamping virtual \parencite{fitzpatrick2017,inkster2018,fulmer2018,guo2025,skjuve2021,maples2024}.

Celah yang tersisa justru terletak pada penyatuannya. Sistem yang cakap menjalankan latihan reflektif tidak otomatis memiliki memori yang dapat ditelusuri lintas sesi, penyampaian skrining suasana hati yang terasa mulus tidak selalu terhubung dengan pengalaman spesifik yang pernah diceritakan pengguna, dan personalisasi yang ditawarkan belum tentu memberi pengguna kendali atas apa yang diingat sistem tentang dirinya. Dengan demikian, kebutuhan yang belum terjawab adalah chatbot pendamping yang menyatukan percakapan reflektif, skrining ringan, memori relasional, dan kendali pengguna dalam satu pengalaman yang tetap sesuai dengan konteks mahasiswa Indonesia.

Tantangan integrasi ini menjadi lebih berat ketika percakapan berlangsung berulang dari waktu ke waktu. Mahasiswa jarang membawa masalah sebagai peristiwa tunggal: keluhan tentang revisi hari ini bisa jadi berkaitan dengan komentar dosen minggu lalu, ekspektasi keluarga, atau rasa tertinggal dari teman seangkatan, sehingga empati yang relevan dan reflektif menuntut sistem memahami mengapa suatu emosi muncul dalam konteks cerita penggunanya \parencite{bickmorePicard2005}. Kegagalan mempertahankan kontinuitas semacam ini terbukti berdampak nyata. \textcite{defreitasReplika2024} mendokumentasikan bahwa perubahan persona Replika setelah pembaruan sistem menimbulkan \textit{identity discontinuity}, kondisi saat pengguna merasa karakter chatbot yang mereka kenal berubah drastis dibandingkan interaksi sebelumnya. \textcite{boine2023} mengaitkan besarnya dampak perubahan semacam itu dengan kuatnya keterikatan emosional yang terbentuk antara pengguna dan pendamping virtualnya. Temuan ini menegaskan bahwa kualitas model bahasa saja tidak menjamin hubungan yang konsisten. Sistem juga harus mampu mempertahankan identitas dan pemahamannya terhadap pengguna dari waktu ke waktu.

Kontinuitas yang semakin kuat juga membawa konsekuensi lain: semakin banyak sistem mengingat, semakin besar pula kebutuhan pengguna untuk tetap merasa memegang kendali. Pada hubungan antarmanusia, seseorang bebas memilih cerita mana yang dibagikan, kapan ingin bicara lebih terbuka, dan kapan ingin menarik diri. Chatbot pendamping yang menyimpan memori jangka panjang perlu menyediakan rasa kendali yang setara, mulai dari mengetahui apa yang diingat, mengoreksi atau menghapus memori tertentu, memilih percakapan yang tidak perlu disimpan, hingga mengakhiri penggunaan tanpa merasa data personalnya terus melekat pada sistem. Tanpa mekanisme semacam ini, memori yang semula dimaksudkan membuat percakapan lebih personal justru berisiko terasa seperti pengawasan.

Persoalan berikutnya terletak pada bentuk memori itu sendiri. Literatur psikologi membedakan \textit{episodic memory}, ingatan atas peristiwa spesifik, dari \textit{semantic memory}, pemahaman yang lebih abstrak dan bertahan lama \parencite{tulving1972}, dan chatbot pendamping idealnya menangkap keduanya. Memori personal juga tidak dapat diperlakukan sebagai arsip statis. Informasi yang dulu penting bisa kehilangan relevansi, pikiran yang dulu diyakini bisa berubah setelah refleksi, dan hubungan yang dulu terasa negatif bisa dimaknai ulang, sehingga memori idealnya memiliki siklus hidup: ditulis ketika relevan, dihubungkan dengan konteks lain, diperbarui saat pemahaman berubah, dan diturunkan kepentingannya ketika tidak lagi sesuai \parencite{zhong2024memorybank}. Pendekatan pencarian berbasis kemiripan semantik yang umum dipakai saat ini berguna untuk menemukan potongan percakapan yang serupa, tetapi memori personal sering menuntut lebih dari kemiripan kata semata. Pertanyaan tentang siapa yang terlibat dalam suatu pengalaman, bagaimana emosi berubah setelah kejadian tertentu, atau pikiran lama mana yang sudah ditafsirkan ulang, tidak selalu terjawab hanya dari potongan teks terdekat. Literatur personalisasi menunjukkan bahwa konteks pengguna yang persisten memang meningkatkan kualitas keluaran model bahasa \parencite{salemi2024lamp}, sementara arah penelitian yang menggabungkan model bahasa dengan graf pengetahuan, termasuk \textit{GraphRAG} dan memori agen berbasis graf temporal, menunjukkan bahwa representasi relasional memberi pegangan yang lebih kuat untuk penalaran lintas entitas, konteks temporal, dan perubahan fakta \parencite{pan2024, edge2024graphrag, rasmussen2025zep, saxenauts2025persona}.

Cara pengguna berinteraksi dengan sistem turut memengaruhi seberapa jauh dukungan ini terasa membantu. Percakapan lisan dapat terasa lebih ringan ketika pengguna belum sanggup merapikan perasaannya menjadi teks \parencite{pradhan2020}, sehingga interaksi suara relevan sebagai fitur pendukung, menurunkan hambatan untuk mulai bercerita, meski tetap membawa konsekuensi latensi karena harus melewati transkripsi, pemrosesan bahasa, dan sintesis suara. Di sisi lain, kedekatan yang terbangun antara chatbot dan pengguna menuntut batas peran yang tetap jelas \parencite{weidinger2022}, terlebih karena pengguna dapat menyampaikan kondisi berat lewat cara yang halus dan tidak langsung \parencite{pendse2022}. Batas keselamatan karena itu perlu hadir sebagai pagar yang menjaga percakapan, bukan tujuan utama yang menggeser fungsi pendampingan.

Pada konteks mahasiswa Indonesia, kebutuhan-kebutuhan tersebut mengambil bentuk yang lebih spesifik. Tekanan akademik dan relasi dengan dosen pembimbing mewarnai masa penyusunan tugas akhir \parencite{pusparini2025thesisstress}. Ekspektasi dan tekanan finansial dari keluarga menambah beban tersendiri \parencite{mardhiyah2026parentalpressure}. Keterlibatan pada organisasi kemahasiswaan berisiko menjadi sumber \textit{burnout} \parencite{octaviani2023organizationalburnout}. Pencarian identitas diri kerap disertai perasaan tidak layak dibanding teman seangkatan \parencite{pakozdy2024impostor}. Ketidakpastian transisi karier membayangi masa menjelang serta sesudah kelulusan \parencite{yang2025employmentanxiety}. Keenam area ini (tekanan akademik, relasi kampus, keluarga, organisasi, identitas diri, dan transisi karier) perlu menjadi acuan domain yang dimiliki oleh mahasiswa. Masalah-masalah tersebut juga banyak diungkapkan melalui bahasa informal dan bahasa campur, dan sistem perlu memahami bahasa itu apa adanya alih-alih memaksanya menjadi bahasa yang lebih baku.

Berangkat dari kesenjangan tersebut, tugas akhir ini mengembangkan AXIS, sebuah chatbot pendamping non-klinis yang dikalibrasi untuk konteks mahasiswa Indonesia. AXIS menyatukan tiga kemampuan inti: percakapan reflektif, asesmen skrining depresi PHQ-9 secara percakapan, dan memori jangka panjang hibrid yang dapat ditulis, dihubungkan, diperbarui, serta dipilih kembali ketika relevan. Ketiga kemampuan tersebut dirancang dalam batas non-klinis dengan kendali tetap berada di tangan pengguna.

Interaksi suara, \textit{Confession Space}, \textit{guardrail}, antarmuka, dan kontrol data diposisikan sebagai kebutuhan pendukung. Fitur-fitur tersebut membantu pengguna mengakses kemampuan inti secara lebih aman dan terkendali, tetapi tidak diperlakukan sebagai kontribusi penelitian yang berdiri sendiri.

Kontribusi rekayasa utama penelitian ini terletak pada dua rancangan memori yang diintegrasikan ke dalam chatbot pendamping. Pertama, skema POLE+O diadaptasi ke domain afektif dengan menambahkan simpul CBT, seperti \textit{Experience}, \textit{Emotion}, \textit{Thought}, \textit{Behavior}, dan \textit{Trigger}, serta relasi \textit{hot cross bun} sebagai struktur graf utama. Kedua, operasi \textit{lifecycle} memori, termasuk \textit{supersession}, \textit{reappraisal}, dan penonaktifan pemicu, dihubungkan dengan perubahan pemaknaan dan label distorsi kognitif. Sepanjang telaah penulis terhadap sistem pada Tabel~\ref{tab:perbandingan-sistem}, kombinasi khusus tersebut belum ditemukan pada \textit{companionship chatbot} yang dikaji. Pernyataan ini menunjukkan posisi kontribusi terhadap karya terpilih yang ditelaah, bukan hasil tinjauan sistematis.

Komponen lain, seperti RRF, MMR, \textit{graph reranker}, \textit{state machine} percakapan, dan penyampaian PHQ-9 secara percakapan, diposisikan sebagai integrasi teknik yang telah digunakan pada domain lain, bukan sebagai kebaruan algoritmik.

Karena tugas akhir ini berada pada tahap pengembangan purwarupa, evaluasi dibatasi pada validitas rekayasa sistem. Bukti primer diperoleh melalui pengujian kontrak fungsional, \textit{benchmark} berlabel, penilaian buta terhadap respons, dan studi ablasi terkendali. Evaluasi pengguna nyata belum dilaksanakan; instrumennya disiapkan sebagai rencana evaluasi lanjutan untuk menilai pengalaman penggunaan. Validasi dampak jangka panjang memerlukan studi longitudinal dengan protokol etik formal yang berada di luar cakupan tugas akhir ini\ifdefined\seminarhasilmode.\else. Keterbatasan ini dinyatakan pada Subbab~\ref{sec:keterbatasan} dan dijabarkan sebagai arah pengembangan pada Subbab~\ref{sec:saran}.\fi
\section{Rumusan Masalah}
\label{sec:rumusan-masalah}

Berdasarkan latar belakang tersebut, tugas akhir ini diarahkan oleh tiga rumusan masalah berikut.

\begin{enumerate}[label=\arabic*.,leftmargin=*,itemsep=4pt]
  \item Bagaimana mengembangkan chatbot pendamping non-klinis yang menghasilkan respons empatik, relevan secara konteks, reflektif, aman, dan sesuai dengan ragam bahasa mahasiswa Indonesia, termasuk bahasa informal dan bahasa campur?

  \item Bagaimana mengintegrasikan mood harian dan asesmen skrining depresi PHQ-9 secara percakapan, berkala, dan sukarela agar pengguna dapat merefleksikan kondisi dua minggu terakhir tanpa pengalaman pendampingan terasa seperti proses diagnosis?

  \item Bagaimana mengembangkan memori jangka panjang yang menjaga kesinambungan percakapan lintas sesi melalui pengenalan dan pemanfaatan konteks permasalahan yang relevan dalam kehidupan mahasiswa Indonesia yang terus berubah?
\end{enumerate}

\section{Tujuan dan Ukuran Keberhasilan Pencapaian}
\label{sec:tujuan}

Tujuan penelitian ini adalah mengembangkan AXIS sebagai chatbot pendamping non-klinis untuk mahasiswa Indonesia. Secara lebih spesifik, penelitian ini bertujuan menghasilkan sistem yang mampu memberikan respons empatik, relevan secara konteks, reflektif, aman, dan sesuai ragam bahasa mahasiswa Indonesia; mengintegrasikan mood harian serta asesmen skrining depresi PHQ-9 secara percakapan; dan mengelola memori jangka panjang yang dapat diperbarui agar konteks pengguna tetap relevan dari waktu ke waktu.

Ukuran keberhasilan ditetapkan pada tingkat klaim sistem secara menyeluruh, bukan pada keberhasilan setiap tugas LLM secara terpisah. Pencapaiannya dinilai sebagai berikut:

\begin{enumerate}[label=\arabic*.,leftmargin=*,itemsep=4pt]
  \item \textbf{Percakapan pendamping.} Respons AXIS dan \textit{baseline} dinilai secara buta pada skenario yang sama untuk menguji kelayakan respons suportif dan melihat nilai tambah konteks memori, bukan untuk menetapkan keunggulan percakapan secara umum. Rubrik ordinal mencakup validasi emosi, pemahaman konteks, eksplorasi, ketepatan teknik CBT, kepatuhan batas non-klinis, kesesuaian ragam bahasa, dan \textit{groundedness} ketika memori digunakan. Pelaporan mencakup distribusi skor per dimensi, preferensi berpasangan, konfigurasi model penilai, dan kesepakatan antarkonfigurasi penilai. Keselamatan dinilai melalui \textit{benchmark} berlabel dengan sensitivitas risiko tinggi sebagai metrik primer, disertai spesifisitas, presisi, $F_2$, jumlah \textit{false negative}, dan kepatuhan respons aman. Hasil bahasa dilaporkan menurut ragam formal, informal, bahasa campur, dan eufemistik.

  \item \textbf{Mood harian dan asesmen skrining depresi.} Fidelity orkestrasi dinilai melalui tingkat kelulusan kontrak untuk penawaran sukarela, penyampaian satu item per giliran, penghormatan penolakan, eskalasi item kesembilan, pemisahan data administratif dari memori, dan penyampaian hasil sebagai skrining. Pemetaan jawaban bebas ke skor 0 sampai 3 dinilai terutama dengan \textit{quadratic weighted kappa}, disertai \textit{macro-F1} dan \textit{accuracy}. Item kesembilan dinilai terpisah melalui sensitivitas, spesifisitas, dan ketepatan rute keselamatan.

  \item \textbf{Memori jangka panjang.} Evaluasi dipisahkan menjadi kemampuan menulis, mengambil, memperbarui, dan menggunakan memori. Penulisan dinilai dengan presisi dan \textit{recall} node serta relasi. Pengambilan dinilai dengan \textit{P@5}, \textit{Recall@5}, MRR, dan nDCG@5. Pembaruan dinilai melalui \textit{update correctness}, sedangkan penggunaan memori dinilai melalui \textit{grounded-answer rate}, \textit{false-personalization rate}, dan probe \textit{stale-belief}. Nilai tambah rancangan hibrid diuji melalui ablasi berpasangan terhadap kondisi vektor saja. \textit{Stale-memory rate}, \textit{contradiction rate}, dan tingkat keberhasilan kueri relasional atau kausal secara umum dicatat sebagai metrik lanjutan yang belum direalisasikan karena belum tersedia korpus berlabel yang memadai; ketiganya tidak digunakan sebagai dasar kesimpulan penelitian.

  \item \textbf{Kendali pengguna.} Kendali pengguna dipertahankan sebagai kebutuhan pendukung lintas rumusan masalah, bukan sebagai klaim penelitian yang berdiri sendiri. Verifikasi pada laporan ini dibatasi pada pengujian fungsional terhadap pengelolaan memori, data personal, dan batas persistensi \textit{Confession Space}. Keberhasilan tugas serta pemahaman pengguna terhadap kontrol tersebut disiapkan untuk evaluasi pengguna nyata pada penelitian lanjutan.
\end{enumerate}

\section{Batasan Masalah}
\label{sec:batasan-masalah}

Agar penelitian tetap fokus, batasan masalah ditetapkan sebagai berikut.

\begin{enumerate}[label=\arabic*.,leftmargin=*,itemsep=4pt]
  \item Sasaran utama sistem adalah mahasiswa Indonesia dewasa pada konteks non-klinis. AXIS tidak ditujukan untuk diagnosis, penentuan terapi, ataupun pengganti psikolog, konselor, psikiater, dan tenaga profesional. Asesmen yang digunakan terbatas pada skrining depresi PHQ-9.

  \item Sistem dikembangkan dan dikalibrasi untuk tekanan akademik, relasi kampus, keluarga, organisasi, identitas diri, dan transisi karier mahasiswa Indonesia. Bahasa yang dicakup adalah bahasa Indonesia formal, informal, dan bahasa campur dengan istilah Inggris yang lazim di lingkungan kampus. Karena evaluasi pengguna nyata belum dilaksanakan, kesesuaian pengalaman AXIS bagi populasi tersebut belum divalidasi melalui partisipan.

  \item Aplikasi web, tampilan mobile, interaksi suara, dan \textit{Confession Space} berfungsi sebagai sarana penggunaan dan kebutuhan pendukung, bukan sebagai kontribusi utama yang dievaluasi secara terpisah. Isi percakapan pada \textit{Confession Space} tidak ditulis ke memori jangka panjang, meskipun batas keselamatan tetap diterapkan.

  \item Evaluasi membuktikan validitas rekayasa purwarupa melalui pengujian kontrak, \textit{benchmark} berlabel, penilaian buta berbasis \textit{LLM-as-judge}, dan ablasi terkendali. Pengujian pengguna nyata belum dilaksanakan. Penelitian ini tidak menilai efektivitas klinis, penurunan gejala depresi, peningkatan kesejahteraan jangka panjang, atau keselamatan klinis penuh.

  \item Keluaran terbuka dinilai secara buta menggunakan \textit{LLM-as-judge} dengan rubrik tertulis, sedangkan kontrak deterministik dinilai terhadap keluaran acuan yang ditetapkan oleh spesifikasi sistem. Penilaian model mengukur kepatuhan terhadap rubrik, bukan pengalaman pengguna.

  \item Untuk kemampuan yang memerlukan cakupan lintas sesi, penelitian ini memakai LongMemEval \parencite{wu2024longmemeval} sebagai kalibrasi eksternal. Untuk kemampuan yang tidak memiliki korpus publik berbahasa Indonesia yang sesuai, seperti penulisan node dan relasi serta pembaruan \textit{lifecycle}, digunakan pengujian skenario-dan-hasil terhadap penulis memori produksi. Teknik tersebut membuktikan perilaku pada kasus yang diuji, bukan estimasi populasi. Metrik penggunaan memori pada jawaban juga diturunkan dari proksi skenario berukuran kecil dan dilaporkan bersama keterbatasannya.

  \item Kendali pengguna dibatasi pada fitur yang tersedia dalam purwarupa, seperti pengelolaan memori dan data personal. Penelitian ini tidak melakukan audit kepatuhan hukum perlindungan data secara menyeluruh.
  \item LLM yang digunakan selama pengembangan diakses melalui penyedia layanan eksternal, yaitu OpenAI dan Google Gemini, tanpa pelatihan model secara mandiri. Oleh karena itu, kebijakan penggunaan, mekanisme keamanan, pembatasan konten, serta perubahan layanan yang ditetapkan oleh masing-masing penyedia berada di luar kendali sistem yang dibangun.

  \item Latensi respons tidak termasuk dalam aspek yang dievaluasi pada penelitian ini. Latensi dapat dipengaruhi oleh berbagai faktor eksternal, seperti perangkat yang digunakan, kondisi jaringan, lokasi server, beban layanan, serta waktu pemrosesan dari \textit{provider} LLM.
\end{enumerate}

\section{Metodologi}
\label{sec:metodologi}

\subsection{Metode Penelitian}
\label{subsec:metode-penelitian}

Metodologi penelitian dalam pengerjaan tugas akhir ini terdiri atas tahap-tahap sebagai berikut.

\begin{enumerate}[label=\arabic*.,leftmargin=*,itemsep=4pt]
  \item \textbf{Identifikasi Permasalahan.} Tahap ini memetakan kebutuhan dukungan emosional non-klinis mahasiswa Indonesia serta menurunkannya menjadi kebutuhan percakapan, mood harian dan asesmen skrining depresi, serta memori jangka panjang, sebagaimana diuraikan pada Subbab~\ref{sec:latar-belakang} dan Subbab~\ref{sec:rumusan-masalah}.

  \item \textbf{Perancangan Solusi.} Tahap ini merancang peran AXIS sebagai chatbot pendamping, alur percakapan reflektif, pemeriksaan suasana hati, memori jangka panjang, dan batas keselamatan yang menjaga sistem tetap non-klinis.

  \item \textbf{Implementasi.} Tahap ini mewujudkan rancangan ke dalam aplikasi web, layanan \textit{backend}, dan penyimpanan data yang mendukung percakapan, mood harian dan asesmen skrining depresi, serta memori jangka panjang.

  \item \textbf{Pengujian dan Evaluasi.} Tahap ini menjalankan validasi kontrak, \textit{benchmark} berlabel, penilaian buta, dan ablasi untuk memeriksa apakah sistem berjalan sesuai rancangan serta apakah konteks memori memberi nilai pada skenario yang diuji. Instrumen pengujian pengguna nyata disiapkan sebagai rencana evaluasi lanjutan, tetapi belum dijalankan pada laporan ini.
\end{enumerate}

\subsection{Metode Pengembangan Perangkat Lunak}
\label{subsec:metode-pengembangan-perangkat-lunak}

Metode pengembangan perangkat lunak yang digunakan adalah iteratif-inkremental. Pengembangan dilakukan modul demi modul secara bertahap: alur percakapan inti dikerjakan dan diuji lebih dahulu, kemudian dilanjutkan dengan lapisan batas keselamatan, mood harian dan asesmen skrining depresi PHQ-9, arsitektur memori jangka panjang hibrid berbasis \textit{knowledge graph}, dan terakhir fitur pendukung seperti interaksi suara. Setiap modul diuji secara fungsional sebelum pengembangan berlanjut ke modul berikutnya, dan rancangan pada modul yang telah selesai tetap dapat direvisi apabila pengujian pada modul selanjutnya mengungkap kebutuhan integrasi baru. Pendekatan ini dipilih karena setiap modul memiliki kebutuhan pengujian yang berbeda. Pengujian batas keselamatan, misalnya, menuntut skenario yang berbeda dari pengujian konsistensi memori, sehingga validasi bertahap per modul lebih sesuai dibandingkan menunda seluruh pengujian ke akhir pengembangan.

\section{Sistematika Pembahasan}
\label{sec:sistematika}

Laporan Tugas Akhir ini membahas pengembangan AXIS sebagai chatbot pendamping non-klinis yang dikalibrasi untuk konteks mahasiswa Indonesia, dengan fokus pada percakapan empatik yang reflektif, mood harian dan asesmen skrining depresi PHQ-9 secara percakapan, serta memori jangka panjang hibrid berbasis \textit{knowledge graph}. \ifdefined\seminarhasilmode Naskah pada tahap seminar hasil ini mencakup empat bab pertama beserta lampiran, sedangkan Bab V disusun menjelang sidang akhir. \else Terdapat lima bab dalam laporan Tugas Akhir ini. \fi

Bab I, yaitu Pendahuluan, memperkenalkan landasan penelitian, menguraikan latar belakang dan konteks permasalahan. Bab ini mencakup rumusan masalah, tujuan dan ukuran keberhasilan pencapaian, serta batasan dari penelitian dan pengembangan, memberikan kerangka kerja yang jelas terhadap hal yang ingin dicapai oleh Tugas Akhir ini. Selain itu, bab ini menjelaskan metode penelitian dan metode pengembangan perangkat lunak yang digunakan.

Bab II, yaitu Kajian Pustaka, berisi tinjauan terhadap literatur yang relevan dan disusun mengikuti urutan tiga rumusan masalah. Bab ini membahas dukungan emosional non-klinis mahasiswa Indonesia, empati dalam sistem percakapan, pendekatan percakapan reflektif, interaksi suara, batas keselamatan, skrining depresi PHQ-9, memori hibrid berbasis graf dan vektor, pengayaan linguistik, serta prinsip arsitektur perangkat lunak.

Bab III, yaitu Deskripsi Solusi, berisi analisis kebutuhan dan rancangan solusi yang menjawab tiga rumusan masalah. Bab ini menjelaskan peran AXIS sebagai chatbot pendamping, rancangan mood harian dan asesmen skrining depresi PHQ-9, rancangan memori jangka panjang, batas keselamatan, dan rancangan evaluasi pada tingkat konseptual.

Bab IV, yaitu Implementasi dan Evaluasi, berfokus pada realisasi teknis dari rancangan pada Bab III dan evaluasi terhadap klaim sistem. Bab ini menyajikan implementasi arsitektur inti, pengujian kontrak fungsional, \textit{benchmark} berlabel untuk keselamatan, PHQ-9, dan memori, penilaian buta respons, studi ablasi \textit{retrieval}, pengendalian keterulangan eksperimen, serta rencana evaluasi lanjutan dengan pengguna nyata.

\ifdefined\seminarhasilmode\else
Bab V, Kesimpulan dan Saran, merupakan bab penutup yang merangkum penemuan dan hasil dari penelitian ini. Bab ini memberikan kesimpulan atas pencapaian tujuan Tugas Akhir, kontribusi yang dapat dipertahankan, keterbatasan sistem yang teridentifikasi pada proses pengembangan dan evaluasi, serta saran untuk pengembangan lanjutan, termasuk arah evaluasi lanjutan dengan persetujuan etik formal dan perluasan cakupan sistem.
\fi
